\section{Future Work}
\label{sec:future-work}

% NOTE: kept your two points, and pulled forward the three extensions your
% own PDF already sketches mathematically in its "Perspectives" section --
% none of which currently appear in main.tex's Future Work outline.

Two directions follow directly from the limitations above. First, where a
mapping between SNOMED~CT and another ontology exists, the same selection and
tokenization mechanism could be applied across ontologies: tokens selected in
one vocabulary could represent concepts from another, extending the
interoperability argument from Section~\ref{sec:introduction} to the token
level itself. \notetoself{This is the same point listed as "(maybe)"
contribution item 5 in the Introduction -- see the note there; keep it in
only one of the two places.} Second, addressing the transductive-selection
limitation directly would mean validating the tokenization on a downstream
clinical task external to the intrinsic metrics of Section~\ref{sec:method:metrics}.

\notetoself{Your existing PDF draft's own "Perspectives" section (Section~8)
already sketches three further extensions, with citations, that don't appear
anywhere in main.tex's Future Work outline. Pulling them in here since they
seem like genuine future-work material rather than something to drop:}

\textbf{Redundancy and token swaps.} The greedy algorithm of
Section~\ref{sec:method:selection} never removes a token once selected. A
token can become nearly redundant once later tokens end up covering the same
concepts it was originally selected for; a marginal-\emph{loss}-based pass
could identify and swap out such tokens after the fact.

\textbf{Primitive vs.\ fully-defined concepts.} Expanding a concept via its
outgoing relationships (Section~\ref{sec:method:tokenizer}) implicitly
assumes that doing so is lossless. For primitive SNOMED~CT concepts -- those
not fully defined by their stated relationships -- this may not hold. A
reconstructibility factor $\alpha(u) \in [0,1]$ ($\alpha(u)=1$ for fully
defined concepts, $\alpha(u)=\eta < 1$ for primitive ones) could discount the
propagated coverage of Eq.~\eqref{eq:coverage} accordingly, while still
treating a primitive concept selected \emph{as its own token} as fully
covered (coverage $1$), since it is then preserved atomically rather than
reconstructed from its relationships.

\textbf{Combining with an ontology-similarity measure.} An independently
developed ontology-based similarity measure \citeneeded{ontology-based
similarity measure referenced in your PDF's Section 8 / reference [6]} could
evaluate how well semantic neighborhoods are preserved and how severe token
collisions are, and could potentially be folded into the marginal-gain
computation itself, e.g.\ as a facility-location-style term -- provided the
submodularity and monotonicity properties the greedy guarantee depends on
(Section~\ref{sec:method:selection}) are re-verified for the modified
objective.
